% Generated draft: requires mathematical and source review.
\documentclass[11pt]{article}
\usepackage[margin=25mm]{geometry}
\usepackage{fontspec}
\setmainfont{Latin Modern Roman}
\usepackage{amsmath,amssymb,mathtools,unicode-math}
\setmathfont{Latin Modern Math}
\usepackage{enumitem,needspace,xurl,hyperref}
\setlength{\parindent}{0pt}
\setlength{\parskip}{5pt}
\setlength{\emergencystretch}{3em}
\newcommand{\N}{\mathbb{N}}
\newcommand{\Z}{\mathbb{Z}}
\newcommand{\Q}{\mathbb{Q}}
\newcommand{\R}{\mathbb{R}}
\newcommand{\C}{\mathbb{C}}
\newcommand{\F}{\mathbb{F}}
\newcommand{\A}{\mathbb{A}}
\newcommand{\PP}{\mathbb{P}}
\newcommand{\E}{\mathbb{E}}
\newcommand{\eps}{\varepsilon}
\newcommand{\dd}{\,\mathrm{d}}
\begin{document}
\section*{UnsolvedMath ID: 701673}
\textbf{Source problem code:} catalogue entry
\subsection{Definitions and mathematical statement}
Let \(\kappa\) be an infinite cardinal. For a set \(X\), let \(\operatorname{Sym}(X)\) denote the group of all bijections from \(X\) to itself, under composition. An embedding of groups means an injective group homomorphism. A group \(G\) is residually finite if, for every \(g\in G\setminus\{1_G\}\), there exist a finite group \(F\) and a homomorphism \(\varphi:G\to F\) such that \(\varphi(g)\ne 1_F\).

\textbf{Question (G. M. Bergman, Kourovka Notebook Problem 20.15).} For every infinite cardinal \(\kappa\) and every residually finite group \(G\), does the existence of a set \(X\) with \(|X|=\kappa\) and an embedding
\[
G\hookrightarrow\operatorname{Sym}(X)
\]
imply that there exist an index set \(I\) with \(|I|=\kappa\), finite groups \(F_i\) for all \(i\in I\), and an embedding
\[
G\hookrightarrow\prod_{i\in I}F_i\,?
\]
Here the product is the unrestricted direct product, with coordinatewise multiplication. The factors need not be distinct, and trivial factors are permitted; thus indexing by exactly \(\kappa\) is equivalent to allowing at most \(\kappa\) factors.

\textbf{Accompanying observations supplied by the source.} The source states that the converse holds: every direct product of \(\kappa\) finite groups is residually finite and embeds in the full permutation group of a set of cardinality \(\kappa\). It also states that the requested implication holds if the number of finite factors is increased to \(2^\kappa\), where \(2^\kappa\) is the cardinality of the power set of a set of cardinality \(\kappa\). The supplied explanation is that an embedding into \(\operatorname{Sym}(X)\) gives \(|G|\leq 2^\kappa\), so residual finiteness supplies at most \(2^\kappa\) homomorphisms to finite groups which together separate every nonidentity element from the identity. These observations are retained as source context, not as a resolution of the question with \(\kappa\) factors.
\subsection{Short English statement}
Must every residually finite group admitting a faithful permutation action on a set of infinite cardinality \(\kappa\) embed into a direct product of \(\kappa\) finite groups? The supplied source states the converse and the version with \(2^\kappa\) factors. The registry labels the question open; that status has not been independently verified.
\subsection{Sources}
\begin{itemize}
\item Local UnsolvedMath registry, permanent ID \(701673\), exact source problem code \texttt{catalogue entry}, title \emph{Kourovka Notebook Problem 20.15}; supplied complete source record, upstream revision \texttt{d371ee917ac9647fbb06cd6fc529deec6f3dda82}. The record was updated on 2026-10-16 and labels the problem \emph{open}; this is supplied registry metadata, not an independently verified literature assessment.
\item \emph{The Kourovka Notebook (Unsolved Problems in Group Theory)}, Problem 20.15, proposed by G. M. Bergman in 2022. The supplied background identifies the 20th issue (2022) and the 21st issue edition (2026), arXiv:1401.0300v48. The supplied collection description names E. I. Khukhro and V. D. Mazurov as editors.
\item G. M. Bergman, \emph{Indag. Math.} \textbf{18} (2007), 349--403, cited in the supplied problem statement. No article title or more specific locator is supplied.
\end{itemize}
\end{document}
